%--------------------------------------------------------------------------------
%
%
%--------------------------------------------------------------------------------
\pagebreak
\section*{WR — Set window radix}
\addcontentsline{toc}{section}{WR — Set window radix}

%--------------------------------------------------------------------------------
\begin{iPageDescEntry}{Syntax}
\texttt{wr [ <radix> [ , <winNum> ]]} \\
\end{iPageDescEntry}

%--------------------------------------------------------------------------------
\begin{iPageDescEntry}{Description}

The \texttt{WR} command sets the numeric display radix for a window.  
The optional \texttt{<radix>} parameter specifies the desired format:
\begin{itemize}
    \item \texttt{HEX} — Display numbers in hexadecimal format.
    \item \texttt{DEC} — Display numbers in decimal format.
\end{itemize}

If the \texttt{winNum} parameter is omitted, the command applies to the current
window. Otherwise, it applies to the specified window.

If the radix parameter is omitted, the command sets the radix to the window default radix.

\end{iPageDescEntry}

%--------------------------------------------------------------------------------
\begin{iPageDescOpEntry}{Example}
\begin{lstlisting}[style=iPageOpStyle]
wn mem, 0x1000      # Create a memory window.
wr HEX              # Display numbers in hexadecimal.
wr DEC, 2           # Window number 2 display in decimal.
wr                  # Set current window radix to default.
\end{lstlisting}
\end{iPageDescOpEntry}

%--------------------------------------------------------------------------------
\begin{iPageDescEntry}{Notes}
Changing the radix affects only the display of numeric data. The underlying values remain unchanged.
For some window types, certain content is always shown in a fixed radix.
For example, addresses in a memory window are always displayed in hexadecimal, while the data values follow the selected radix.
\end{iPageDescEntry}
